\section{Appendix I: Variations on Abhyankar's lemma}
\label{XIII.5}
%
This appendix contains various variants of Abhyankar's lemma.

\begin{proposition}
\label{XIII.5.1} Let $X=\Spec A$ be a regular local scheme,
$D=\sum_{1\leq i\leq r} \divisor f_i$ a divisor with normal
crossings, where the $f_i$ are elements of the
maximal ideal of~$A$ forming part of a regular system of
parameters. Let $n_i$, $1\leq i\leq r$, be integers $\geq 0$
and set
$$
X'=X[T_1,\dots,T_r]/(T_1^{n_1}-f_1,\dots,T_r^{n_r}-f_r),
$$
$U'=U\times_X X'$. Then $X'$ is regular and $U'$ is the
complement in $X'$ of the divisor with normal crossings
$\sum_{1\leq i\leq r}\divisor T_i$. If the integers $n_i$ are prime
% original p. 428
to the residue characteristic $p$ of~$X$, $U'$ is a
connected \'etale covering of~$U$, tamely
ramified relative to $D$ \textup{(\Ref{XIII.2.3}~c))}.
\end{proposition}

Indeed $X'$ is the spectrum of a local ring $A'$ whose
maximal ideal is generated by $T_1,\dots,T_r$.
(We may assume $r=\dim(A)$.)
Since $A'$ is finite and flat over~$A$, hence of dimension $r$, $A'$ is
regular (EGA~$0_{\textup{IV}}$~17.1.1) and the $T_i$ form a regular
system of parameters of~$A'$. Suppose the $n_i$ prime
to $p$. Since all the $f_i$ are invertible on~$U$, the fact that
$U'$ is \'etale over~$U$ follows from I~\Ref{I.7.4}. Moreover $U'$
is tamely ramified relative to $D$; indeed, let
$x_i$ be the generic point of~$V(f_i)$, $\bar{R}$ the
strict localization of~$R=\cal{O}_{X,x_i}$, $\bar{K}$ the field of
fractions of~$\bar{R}$. Then the $\bar{K}$-algebra that
represents $U'|\bar{K}$ is obtained from the field
$\bar{K}[T_i]/(T_i^{n_i}-f_i)$ by making an unramified
extension; it is therefore tamely ramified over~$\bar{R}$.

\begin{proposition}[Absolute Abhyankar lemma]
\label{XIII.5.2}
%
Let $X$ be a regular local scheme,
$$
D=\sum_{1\leq i\leq r} \divisor f_i
$$
a divisor with normal crossings as in~\Ref{XIII.5.1},
$Y=\Supp D$, $U=X-Y$. Let~$V$ be an \'etale covering of~$U$
tamely ramified relative to $D$. If $x_i$ is the
generic point of the closed set $V(f_i)$, $\cal{O}_{X,x_i}$ is
a discrete valuation ring with field of fractions $K_i$, and
we have $V|K_i=\Spec(\prod_{j\in J_i}L_j)$, where the $L_j$ are
finite separable extensions of~$K_i$; denote by $n_j$ the order of the
inertia group of a Galois extension generated by $L_j$
and by $n_i$ the l.c.m.\ of the $n_j$ as $j$ runs through $J_i$. If we set
$$
X'=X[T_1,\dots,T_r]/(T_1^{n_1}-f_1,\dots,T_r^{n_r}-f_r),
$$
$U'=U_{(X')}$, $V'=V_{(X')}$, etc., the \'etale covering $V'$
of~$U'$ extends in a unique way, up to unique isomorphism,
to an \'etale covering of~$X'$, and the $n_i$ are
prime to the residue characteristic $p$ of~$X$.
\end{proposition}

Uniqueness follows from the fact that $X'$ is normal
\eqref{XIII.5.1}; indeed
% original p. 429
an \'etale covering of~$X'$ that extends $V'$ is isomorphic to the
normalization of~$X'$ in the fiber of~$V'$ at the generic point
of~$X'$ \eqref{I.10.2}. If $\bar{x}'$ is a geometric point
of~$Y'$, we denote by $\bar{X}'$ the strict localization of~$X'$ at
$\bar{x}'$, $\bar{V}'=V'_{(\bar{X}')}$, etc. By descent, taking
uniqueness into account, it suffices to show that, for every
geometric point $\bar{x}'$ of~$Y'$, the \'etale covering
$\bar{V}'$ of~$\bar{U}'$ extends to $\bar{X}'$. Given
that an \'etale covering of an open subset of the regular scheme~$X'$
which contains all the points $x'$ such that $\dim
\cal{O}_{X',x'}\leq 1$ extends to all of $X$
(SGA~2 XIV~1.11), we may even restrict ourselves to the points $\bar{x}'$ which
project onto a maximal point of~$Y'$. Now, at such a point
$\bar{x}'$, the fact that $\bar{V}'$ extends to an \'etale
covering of~$\bar{X}'$ follows from~\Ref{X.3.6}.

Let us show that the $n_i$ are prime to $p$. Indeed, if this
were not so, we would have for example $p|n_1$. Replacing
$X$ by
$$X[T_1,\dots,T_r]/(T_1^{n_1/p}-f_1,T_2^{n_2}-f_2,\dots,T_r^{n_r}-f_r),$$
we reduce to the case where $X'=X[T_1]/(T_1^p-f_1)$. It suffices
to show that $V$ extends to an \'etale covering of~$X$,
because we will then have $n_1=1$, contrary to the hypothesis. For this we may
assume that $X$ is strictly local. Let $Z$ be the
closed subscheme of~$X$ with equation $p=0$ and $Z_1=Z\cap
X_{f_1}$; $Z_1$ is a nonempty open subset of~$Z$. By what
precedes, the \'etale covering $V'$ of~$U'$ extends
to an \'etale covering $W'$ of~$X'$. Let $W''_1$ and $W''_2$
be the inverse images of~$W'$ under the two projections $X''=X'\times_X
X'\rightrightarrows X'$, and let us show that the descent isomorphism
$u\colon W''_1|U''\to W''_2|U''$ extends to an $X''$-morphism
$W''_1\to W''_2$, which will necessarily be a descent
datum on~$W'$ relative to $X'\to X$. Let $Z''$
(\resp $Z''_1$) be the inverse image of~$Z$ (\resp $Z_1$) in $X''$. Since
the morphism $Z''\to Z$ is radicial, there exists an isomorphism
$v\colon W''_1|Z''\to W''_2|Z''$ which extends the isomorphism
$u|Z''_1$. But, since $X$ is henselian, we have a bijection
$$
\Hom_{X''}(W''_1,W''_2)\simeq\Hom_{Z''}(W''_1|Z'',W_2''|Z''),
$$
whence a morphism $w\colon W''_1\to W''_2$ lifting $v$. The
subscheme, at once open
% original p. 430
and closed, of~$X''$ over which $u$ and $w$ coincide
contains $Z''_1$, hence is equal to~$X''$, whence the fact that $V$
extends to~$X$.

\setcounter{subsection}{3} \setcounter{subsubsection}{-1}

\subsubsection{}
\label{XIII.5.3.0} Let us take up again the hypotheses and the notation
of~\Ref{XIII.5.2}, assuming moreover
$X$
strictly local. Then
it follows from \loccit that every connected \'etale covering
of~$U$ tamely ramified relative to $D$ is a
quotient of a tamely ramified covering of the
form
$$
U'=U[T_1,\dots,T_r]/(T_1^{n_1}-f_1,\dots,T_r^{n_r}-f_r),
$$
where the $n_i$ are integers prime to $p$. Let
$\bbmu_n$ be the group of $n$th roots of unity of
$U$. The group of $U$-automorphisms of~$U'$ is none other than the
group $\bbmu_{n_1}\times\dots\times\bbmu_{n_r}$, an
$n_i$th root of unity $\xi_i$ acting on~$U'$ by
transforming $T_i$ into $\xi_i T_i$. We therefore have the following statement:

\setcounter{subsection}{2}

\begin{corollary}
\label{XIII.5.3} Let $X$ be a regular strictly local scheme
of residue characteristic $p\geq 0$, $D=\sum_{1\leq i\leq
n} \divisor f_i$ a divisor with normal crossings on~$X$,
$U=X-\Supp D$. Set
$$
\tilde U=\varprojlim_{(n_i)}
U[T_1,\dots,T_r]/(T_1^{n_1}-f_1,\dots,T_r^{n_r}-f_r),
$$
the projective limit being taken over the filtered ordered
set (for the divisibility relation) of families of integers
$n_i>0$, prime to $p$. Then $\tilde U$ is a tamely ramified
universal covering of~$U$. Consequently the tamely ramified
fundamental group of~$U$ is
$$
\pi_1^{\tame}(U)\simeq \prod_{\ell\neq p} \ZZ_\ell [1]^r
\qquad \text{(canonical isomorphism)}
$$
where we have set
$\ZZ_\ell[1]=\varprojlim_{n>0}\bbmu_{\ell^n}$. The group
$\pi_1^{\tame}(U)$ is non-canonically isomorphic to $\prod_{\ell\neq p}
\ZZ_\ell^r$.
%
\end{corollary}

\begin{proposition}
\label{XIII.5.4} Let $f\colon X\to S$ be a morphism of
schemes, $D=\sum_{1\leq i\leq r} \divisor f_i$ a divisor with
normal crossings relative to~$S$ \eqref{XIII.2.1}, where,
% original p. 431
for each point $x$ of~$Y=\Supp D$, if $I(x)\subset[1,r]$ is
the set of $i$ such that $f_i(x)=0$, the subscheme
$V((f_i)_{i\in I(x)})$ is smooth over~$S$ of codimension
$\card I(x)$
in $X$. Let $U=X-Y$. Let $x$ be a point of
$Y$, $X_1=\Spec\cal{O}_{X,x}$, $U_1=U\times_X X_1$, $n_i$, $i\in I(x)$,
integers, and
$$
X'=X[T_i]_{i\in I(x)}/(T_i^{n_i}-f_i).
$$
Then, if $x'$ is the point of~$X'$ over~$x$, $X'$ is smooth
over~$S$ at $x'$. If the integers~$n_i$ are prime to the
characteristic $p$ of~$\kres(x)$, $U'_1=U_1\times_X X'$ is a
connected \'etale covering of~$U_1$ tamely
ramified over~$X_1$ relative to~$S_1$ \eqref{XIII.2.1.1}.
\end{proposition}

If $s=f(x)$, the geometric fiber $X'_{\bar{s}}$ is
regular at the point $x'$ \eqref{XIII.5.1}; since $X'$ is flat
over~$S$ in a neighborhood of~$Y$, this proves that $X'$ is smooth over~$S$ at
$x'$ (EGA~IV~12.1.6). If the integers $n_i$ are prime to $p$,
$U'_1$ is an \'etale covering of~$U_1$ \eqref{I.7.4}; it is
tamely ramified over~$X_1$ relative to~$S$ because
this is so on the geometric fibers at each point of~$S$
\eqref{XIII.5.1}. Finally the fact that $U'_1$ is connected follows from
(SGA~4 XVI~3.2).

\begin{proposition}[Relative Abhyankar lemma]
\label{XIII.5.5}
Let $X$ be an $S$-scheme, $D$ a
divisor with normal crossings relative to~$S$, as
in~\eqref{XIII.5.4}. Let $Y=X-\Supp D$,
$U=X- Y$, $x$ a point of~$Y$, $X_{1}$ the strict
localization of~$X$ at a geometric point over~$x$,
$U_{1}=U\times_{X}X_{1}$, $V_{1}$ an \'etale covering of
$U_{1}$. Assume that, for every maximal point $s$ of~$S$,
$V_{1\overline{s}}$ is tamely ramified over~$X_{1\overline{s}}$ relative to $\overline{s}$. Then one can
find integers $n_{i}$ prime to the characteristic $p$
of~$\kres(x)$, with $i\in I(x)$, such that, if we set
$$
X'_{1}=X_{1}[T_{i}]_{i\in I(x)}\big/ (T_{i}^{n_{i}}-f_{i}) \quoi,
$$
$U'_{1}=U_{1}\times_{X_{1}}X'_{1}$, etc., the \'etale covering
$V'_{1}$ of~$U'_{1}$ extends in a unique way, up to
unique isomorphism, to an \'etale covering of
$X'_{1}$. In particular
% original p. 432
$V_{1}$ is tamely ramified over~$X_{1}$ relative
to~$S$.
\end{proposition}

We may assume $S$ local noetherian with closed point~$f(x)$.
For each maximal point $s$ of~$S$ and for each $i\in I(x)$, let
$x_{i}$ be the generic point of the closed set $V(f_{i})$ of the fiber
$X_{1\overline{s}}$. The local ring $(\cal{O}_{X_{1},x_{i}})_{\red}$
is a discrete valuation ring with field of fractions $K_{i}$
and we have
$V_{1}|K_{i}=\Spec(\prod_{j\in I(x_{i})}L_{j})$,
where $L_{j}$ is a finite separable extension of~$K$; we denote by
$n_{j}$ the order of the inertia group of a Galois extension
generated by $L_{j}$ and by $n_{i}$ the l.c.m.\ of the $n_{j}$ as
$s$ runs through the maximal points of~$S$ and
$j\in I(x_{i})$.

The $n_{i}$ being thus chosen, we shall show that $V'_{1}$
extends in a unique way to an \'etale covering of
$X'_{1}$. Uniqueness follows from the fact that, $X'$ being smooth
over~$S$ at the points of~$Y$, we have $\prof\et_{Y'_{1}}(X'_{1})\geqslant 2$
(SGA~4 XVI~3.2 or SGA~2 XIV~1.19). Let $x'_{1}$ be a point of
$Y'_{1}$, $\overline{x}'_{1}$ a geometric point over~$x'_{1}$,
and denote by $\overline{X'_{1}}$ the strict localization of
$X'_{1}$ at $\overline{x}'_{1}$,
$\overline{U'_{1}}=U'_{1(\overline{X'_{1}})}$,~etc. By descent,
taking uniqueness into account, it suffices to show that
$\overline{V}'_{1}$ extends to $\overline{X'_{1}}$. Moreover we
may restrict ourselves to taking for $x'_{1}$ the maximal points of
$Y'_{1}$; indeed we will then have an extension of~$V'_{1}$ over an
open subset $W'_{1}$ of~$X'_{1}$ containing the maximal points of~$Y'_{1}$;
now, if $Z'_{1}=X'_{1}- W'_{1}$, we have
$\codim(Z'_{1s},X'_{1s})\geqslant 2$ if $s$ is a maximal point of
$S$ and $\codim(Z'_{1s},X'_{1s})\geqslant 1$, $\prof\et_{s}(S)\geqslant
1$ if $s$ is a point of~$S$ which is not maximal; the fact that
$V'_{1}$ extends to all of $X'_{1}$ then follows
from~(SGA~2 XIV~1.20). But, at a geometric point
$\overline{x}'_{1}$ over a maximal point of~$Y'_{1}$,
$X'_{1\red}$ is the spectrum of a discrete valuation ring, and
the fact that $\overline{V'_{1}}$ extends to $\overline{X'_{1}}$
then follows from~X~\Ref{X.3.6}.

Let us show that the $n_{i}$ are prime to $p$. Indeed, if this
were not so, there would be an index $i_{0}\in I(x)$ such that $p$
divides $n_{i_{0}}$. Replacing $X$ by
$X_{1}[T_{i_{0}},T_{i}]_{i\in I(x)}\big/
(T_{i}^{n_{i_{0}}/p}-f_{i_{0}},T_{i}^{n_{i}}-f_{i})$, we reduce
to the case where $X'_{1}=X_{1}[T]/T^{p}-f_{i_{0}}$. By what
precedes, the \'etale
% original p. 433
covering $V'_{1}$ of~$U'_{1}$ extends to an \'etale
covering $E'_{1}$ of~$X'_{1}$. Let $\eta$ be the closed point of
$S$; since the morphism $X'_{1\eta}\to X_{1\eta}$ is radicial,
$V_{1\eta}$ extends to an \'etale covering $E_{1\eta}$ of
$X_{1\eta}$. One then deduces, as in~\Ref{XIII.5.2}, that
$E'_{1}$ is endowed with a descent datum relative to the
morphism $X'_{1}\to X_{1}$, which extends the natural descent
datum that we have on~$E'_{1}|U'_{1}$. It follows that
$V_{1}$ extends to $X_{1}$; but this implies that we have
$n_{i_{0}}=1$, contrary to the hypothesis $n_{i_{0}}=p$.

\begin{corollary}
\label{XIII.5.6}
Let $X$ be an $S$-scheme, $D=\sum_{1\leqslant i\leqslant
r}\divisor f_{i}$ a divisor with normal crossings relative
to~$S$, as in~\Ref{XIII.5.4}. Let $\overline{x}$ be a geometric
point of~$X$, $\overline{X}$ the strict localization of
$X$ at $\overline{x}$, $\overline{Y}=Y_{(\overline{X})}$,
$\overline{U}=\overline{X}-\overline{Y}$ and
$$
\widetilde{U}= \varprojlim_{(n_{i})}
\overline{U}[T_{i}]_{i\in I(x)}\big/ (T_{i}^{n_{i}}-f_{i}) \quoi,
$$
the projective limit being taken over the filtered set of
families of integers $n_{i}>0$, prime to the characteristic
$p$ of~$\kres(x)$. Then $\widetilde{U}$ is a tamely ramified
universal covering of~$\overline{U}$ relative to~$S$. Consequently
the tamely ramified fundamental group
of~$\overline{U}$ is
$$
\pi_{1}^{\tame}(\overline{U})\simeq \prod_{\ell\neq
p}\ZZ_{\ell}[1]^{I(x)} \qquad \textrm{(canonical isomorphism)\quoi.}
$$
The group $\pi_{1}^{\tame}(\overline{U})$ is non-canonically
isomorphic to $\prod_{\ell\neq p}\ZZ_{\ell}^{I(x)}$.
\end{corollary}

\begin{subremark}
\label{XIII.5.6.1}
Let $X$ be an $S$-scheme, $D=\sum_{1\leqslant i\leqslant
r}\divisor f_{i}$ a divisor with normal crossings relative
to~$S$, as in~\Ref{XIII.5.4}, $U=X-\Supp D$.
For every subset $I\subset[1,r]$ let
$$
X_{I}= \Big(\tbigcap_{i\in I}V(f_{i})\Big) \cap \Big(\tbigcap_{i\in\complement
I}X_{f_{i}}\Big)\quoi .
$$
Let $p$ be an integer that is prime or zero and let $Z$ be a subset of
$X_{I}$ all of whose points are of characteristic~$p$. Let
$$
\widetilde{U}_{I}= \varprojlim_{(n_{i})}
U[T_{i}]_{i\in I}
\big/ (T_{i}^{n_{i}}-f_{i}) \quoi,
$$
where
% original p. 434
the projective limit is taken over the filtered set of
families of integers $n_{i}>0$, prime to $p$. Then, for every
geometric point $\overline{x}$ of~$Z$, the inverse image of
$\widetilde{U}_{I}$ on~$\overline{U}$ is identified with the tamely
ramified universal covering of~$\overline{U}$.
\end{subremark}

\begin{corollary}
\label{XIII.5.7}
The notation is that of~\Ref{XIII.5.6}. Let $\overline{S}$ be the
strict localization of~$S$ at~$\overline{x}$, and
$$
\overline{g}\colon\overline{U}\to\overline{S} \quad \text{and} \quad
\tilde{g}\colon\widetilde{U}\to\widetilde{S}
$$
the canonical morphisms. Then the morphisms $\overline{g}$ and
$\tilde g$ are $0$-acyclic \textup{(SGA~4 XV~1.3)}. Let $G$ be a
constructible sheaf of groups on~$\overline{S}$, $F=\overline{g}^*G$, $P$
a torsor under $F$. Then, for $P$ to be tamely
ramified over~$\overline{X}$ relative to $\overline{S}$, it
is necessary and sufficient that its inverse image $\widetilde{P}$ on~$\widetilde{U}$ be trivial.
\end{corollary}

Indeed, for every scheme
$\overline{X'}=\overline{X}[T_{i}]_{i\in I(x)}\big/
(T_{i}^{n_{i}}-f_{i})$, where the $n_{i}$ are integers~$>0$
prime to $p$, the morphism
$\overline{f}'\colon\overline{X'}\to\overline{S}$ is $0$-acyclic.
The geometric fibers of~$\overline{f}'$ at the various
points of~$\overline{S}$ are therefore connected and even
irreducible. The same therefore holds for the geometric fibers
of the morphisms
$\overline{g}'\colon\overline{U'}\to\overline{S}$, which proves that
the $\overline{g}'$, hence also $\tilde g$, are
$0$-acyclic (SGA~4 XV~1.16).

It is clear that a torsor $P$ on~$\overline{U}$ with group $F$, whose
inverse image on~$\widetilde{U}$ is trivial, is
tamely ramified over~$\overline{X}$ relative to
$\overline{S}$. Let us show that conversely, if $P$ is
tamely ramified over~$\overline{X}$ relative to
$\overline{S}$, its inverse image on~$\widetilde{U}$ is trivial.

It follows from (SGA~4 IX~2.14~(ii)) that one can find a
finite morphism $n\colon S_{1}\to\overline{S}$, a constant sheaf of groups
$C$ on~$S_{1}$, a monomorphism $G\to n_*C$.
Consider the following commutative diagram made up of
cartesian squares:
$$
\xymatrix{ \widetilde{U}_{1} \ar[r] \ar[d]_{r} &U_{1} \ar[r]^{g_{1}}
\ar[d]_{q} &S_{1} \ar[d]_{n} \cr \widetilde{U} \ar[r] &\overline{U}
\ar[r]^{\overline{g}} &\overline{S} \rlap{\quoi.}\cr }
$$
Let
% original p. 435
$C_{1}$ (\resp $\widetilde{C}_{1}$) be the inverse image of~$C$ on~$U_{1}$
(\resp on~$\widetilde{U}_{1}$). We have a commutative diagram, in
which $i$ and $j$ are isomorphisms (SGA~4 VIII~5.8):
\begin{equation*}
\label{eq:XIII.5.7.*}
\tag{$*$}
\begin{array}{c}
\xymatrix{ \H^{1}(\overline{U},q_*C_{1}) \ar[r]^{i}
\ar[d] &\H^{1}(U_{1},C_{1}) \ar[d] \cr
\H^{1}(\widetilde{U},r_*\widetilde{C}_{1}) \ar[r]^{j}
&\H^{1}(\widetilde{U}_{1},\widetilde{C}_{1}) \rlap{\quoi.}\cr }
\end{array}
\end{equation*}
Let $Q$ be the torsor under $q_*C_{1}$ deduced from~$P$ by
the extension of the structure group $F\to q_*C_{1}$.
By~\Ref{XIII.2.1.4}, $Q$ is tamely ramified
over~$\overline{X}$ relative to $\overline{S}$. To the torsor $Q$
there corresponds, by means of $i$, a torsor $Q_{1}$ under $C_{1}$, and
it is clear that $Q_{1}$ is tamely ramified over~$X_{1}=\overline{X}\times_{\overline{S}}S_{1}$ relative to~$S_{1}$. It therefore follows from~\Ref{XIII.5.6} that the inverse image
$\widetilde{Q}_{1}$ of~$Q_{1}$ on~$\widetilde{U}_{1}$ is trivial,
and diagram~\eqref{eq:XIII.5.7.*} then shows that the inverse image
$\widetilde{Q}$ of~$Q$ on~$\widetilde{U}$ is trivial.

Consider the following commutative diagram, whose second
row is exact (SGA~4 XII~3.1):
\begin{equation*}
\label{eq:XIII.5.7.**}
\tag{$**$}
\begin{array}{c}
\xymatrix{ \H^{0}(\overline{S},n_*C/G) \ar[r]
\ar[d]_{k} &\H^{1}(\overline{S},G)\rlap{$~=~1$} \ar[d] \cr
\H^{0}(\widetilde{U},r_*\widetilde{C}_{1}/\widetilde{F}) \ar[r]
&\H^{1}(\widetilde{U},\widetilde{F}) \ar[r]
&\H^{1}(\widetilde{U},r_*\widetilde{C}_{1}) \rlap{\quoi.} \cr }
\end{array}
\end{equation*}
Since the morphism $\widetilde{U}\to\overline{S}$ is $0$-acyclic,
$k$ is an isomorphism. The fact that $\widetilde{P}$ is trivial
then follows from~\eqref{eq:XIII.5.7.**}.

\section[Appendix II: finiteness for direct
images of stacks]{Appendix II: finiteness theorem for direct
images of stacks}
\label{XIII.6}

\begin{proposition}
\label{XIII.6.1}
Let $S$ be a locally noetherian scheme, $f\colon X\to S$
a morphism. If $S'$ is an $S$-scheme, we denote by $X'$
(\resp $f'$, etc.) the inverse image of~$X$ (\resp $f$, etc.) under the
morphism $S'\to S$. Suppose that, for every
% original p. 436
scheme $S'$ \'etale over~$S$ and every constructible sheaf of sets
$F$ on~$X'$, $f'_*F$ is constructible, and that, for
every constructible sheaf of groups $F$ on~$X'$, $\R^{1}f'_*F$
is constructible. Let $\Phi$ be a $1$-constructible stack on~$X$~\eqref{XIII.0}. Then $f_*\Phi$ is $1$-constructible.
\end{proposition}

For every scheme $S'$ \'etale over~$S$ and for every object $x$ of
$(f_*\Phi)_{S'}$, we have an isomorphism
$$
\SheafAut_{S'}(x)\simeq f'_*\SheafAut_{X'}(x) \quoi,
$$
where, in the right-hand side of the equality, $x$ is
regarded as an object of~$\Phi_{X'}$. The hypotheses
made therefore imply that $f_*\Phi$ is constructible.
Let $S\Phi$ be the sheaf of maximal subgerbes of
$\Phi$~\cite[III~2.1.7]{XIII.1}. Since $f_*(S\Phi)$ is constructible, we
can apply (SGA~4 IX~2.7) to it, and the fact that $f_*\Phi$ is
$1$-constructible then follows from the lemma which follows.

\begin{sublemma}
\label{XIII.6.1.1}
Let $S$ be a locally noetherian scheme, $f\colon X\to S$
a morphism, $\Phi$ a stack on~$X$. Suppose given a sheaf
on~$S$, representable by an \'etale $S$-scheme of finite
type $T$, a surjective morphism
$$
a\colon T\to f_*(S\Phi )
$$
and an object $p$ of the fiber $\Phi_{X_T}$ (where $X_T=X\times_ST$),
defining in $f_*(S\Phi )(T)=\allowbreak S\Phi (X_T)$ an element
equal to the image $q$ under $a$ of the identity section of
$T(T)$. Let $f_T\colon X_T\to T$ be the canonical morphism and suppose
that the sheaf $\R^1f_{T*}(\SheafAut_{X_T}(p))$ is constructible;
then so is~$S(f_*\Phi )$.
\end{sublemma}

The canonical morphism $f^*f_*\Phi\to\Phi$ gives a morphism
$$
S(f^*f_*\Phi )\simeq f^*(S(f_*\Phi ))\to S\Phi ,
$$
whence a canonical morphism
$$
\varphi\colon S(f_*\Phi )\to f_*(S\Phi ).
$$
Let
% original p. 437
$F=S(f_*\Phi )$ and let $G$ be the image of~$F$ under $\varphi$; by
(SGA~4 IX,~2.9) $G$ is a constructible sheaf.

It suffices to show that, for every point $s$ of~$S$, there exists a
nonempty open set~$U$ of~$s$ such that $F|U$ is locally constant
constructible. Let $s\in S$, $\overline{s}$ a geometric
point over~$s$, $\overline{q}_1,\dots,\overline{q}_n$ the elements of~$G_{\overline{s}}$. By
definition of~$T$, there exist $S$-morphisms
$h_i\colon\overline{s}\to S'$ such that
$\overline{q}_i=h_i^*(q)$. Let $S'$ be the fiber product over~$S$ of
$n$ schemes isomorphic to $T$, $\overline{s}\to S'$ the fiber
product of the $h_i$, $X'=X\times_SS'$, $q_i$ (\resp $p_i$) the inverse
image of~$q$ (\resp $p$) under the $i$th projection of~$S'$ onto~$T$. If $\Psi_i$ is the maximal subgerbe of~$\Phi|X'$
generated by $p_i$, the sheaf
$F_i=\R^1f_*'(\SheafAut_{X'}(p_i))$ is none other than the sheaf
$S(f_*'\Psi_i)$ of maximal subgerbes of~$f_*'\Psi_i$. In
particular the canonical injection $\Psi_i\to\Phi |X'$ gives a
morphism
$$
\alpha_i\colon F_i\to F|S'.
$$

We shall show that $\alpha_i$ is a bijection of~$F_i$ onto
the inverse image of~$q_i$ in $F|S'$. For every scheme $S''$
\'etale over~$S'$, every section $y$ of~$F_i(S'')$ has as image
$q_i|S''$ in $F(S'')$, because, locally for the \'etale topology
on~$S''$, $y$ is defined by an object~$x$ of~$\Phi_{X''}$ which is
isomorphic to $p_i|X''$. Conversely, if $y\in F(S'')$ has as
image $q_i|S''$ in $F(S'')$, then locally for the \'etale topology
on~$S''$, $y$ is defined by an object~$x$ of~$\Phi_{X''}$ which is
isomorphic to $p_i$; consequently $x$ is an object of~$\Psi_{iX''}$ and
consequently $y\in F_i(S'')$.

The proof is completed by using~\Ref{XIII.6.1.2}
below. Indeed, one can find an open neighborhood $U'$ of~$s$
such that $q_1|U',\dots,q_n|U'$ are sections of~$G(U')$ and
generate this sheaf. Since the $F_i|U'$ and $G|U'$ are
constructible, so is~$F|U'$,
by~\Ref{XIII.6.1.2}; replacing $U$ by a smaller
open set, $F|U$ is locally constant, which completes the
proof.

\begin{sublemma}
\label{XIII.6.1.2}
Let
% original p. 438
$S$ be a locally noetherian scheme, $F\to G$ a surjective
morphism of sheaves of groups on~$S$. Let $q_i$ be a
finite family of sections of~$G$ over~$X$ which generate $G$, and, for
each $i$, let $F_i$ be the subsheaf of~$F$ which is the inverse image of
$q_i$. Then, if $G$ and the $F_i$ are constructible, so is~$F$.
\end{sublemma}

To prove that $F$ is constructible, it suffices to show that, for
every point~$s$ of~$S$, there exists an open neighborhood $U$ of~$s$ such
that $F|U$ is locally constant constructible. So let $s$ be a
point of~$S$. Since the sheaves $F_i$ and $G$ are constructible, one
can find an open neighborhood $U$ of~$s$ such that $F_i|U$ and $G|U$
are locally constant. Let us then show that $F|U$ is locally
constant. By (SGA~4 IX~2.13~(i)), it suffices to see that, if
$\overline{s}$ is a geometric point over~$s$,
$\tilde{s}$ a geometric point of~$U$ and
$\overline{s}\to\tilde{s}$ a specialization morphism, the
canonical morphism
$$
F_{\tilde{s}}\to F_{\overline{s}}
$$
is bijective.

We consider the commutative diagrams
$$
\xymatrix{ (F_i)_{\tilde{s}} \ar[r]^-{\tilde{a}_i} \ar[d]^{\wr} &
F_{\tilde{s}} \ar[r]^-{\tilde{a}} \ar[d]_b & G_{\tilde{s}}
\ar[d]^{\wr}\\ (F_i)_{\overline{s}} \ar[r]^-{\overline{a}_i} &
F_{\overline{s}} \ar[r]^-{\overline{a}} & G_{\overline{s}} }
$$
Let $\overline{q}_i$ (\resp $\tilde{q}_i$) be the inverse image of~$q_i$
in $G_{\overline{s}}$ (\resp $G_{\tilde{s}}$); the morphisms
$\overline{a}$ and $\tilde{a}$ are surjective, and the morphism
$\overline{a}_i$ (\resp $\tilde{a}_i$) induces a bijection of
$(F_i)_{\overline{s}}$ (\resp $(F_i)_{\tilde{s}}$) onto~$\overline{a}^{-1}(q_i)$ (\resp onto~$\tilde{a}^{-1}(q_i)$). It therefore follows
from the diagram above that $b$ is an isomorphism.

\begin{corollary}
\label{XIII.6.2}
Let
% original p. 439
$S$ be a locally noetherian scheme, $f\colon X\to S$
a proper morphism. Let $\Phi$ be a $1$-constructible stack on~$X$;
then $f_*\Phi$ is a \hbox{$1$-constructible} stack.
\end{corollary}

The proof of~\Ref{XIII.6.1} also proves the following
result, taking \Ref{XIII.2.4}~2) into account.

\begin{corollary}
\label{XIII.6.3}
Let $S$ be a locally noetherian scheme, $f\colon X\to S$ a
morphism, $D$ a divisor on~$X$ with normal crossings
relative to~$S$ \eqref{XIII.2.1}, $Y=\Supp D$, $U=X-Y$,
$i\colon U\to X$ the canonical immersion. Let $\Phi$ be a stack on~$U$
given, locally for the \'etale topology on~$X$ and $S$,
as the inverse image of a $1$-constructible stack $\Psi$ on~$S$. Then
the stack $i_*^{\tame}\Phi$ is $1$-constructible.
\end{corollary}

\begin{thebibliography}{0}{XIII.7}

\bibitem[1]{XIII.1} J\ptbl \textsc{Giraud}, Cohomologie non ab\'elienne de
degr\'e~$2$, thesis, Paris (1966).

\bibitem[2]{XIII.2} J\ptbl \textsc{Giraud},
Cohomologie non ab\'elienne, Springer, Berlin-New\kern2pt York (1971).

\bibitem[3]{XIII.3} 
H\ptbl\textsc{Seifert}, W\ptbl\textsc{Threlfall}, Lehrbuch der Topologie, Chelsea, New
York (1934).


\bibitem[4]{XIII.4}
J-P\ptbl\textsc{Serre}, Cohomologie galoisienne, Lecture Notes
\no 5, Springer, Berlin (1964).

\bibitem[5]{XIII.5}
J-P\ptbl\textsc{Serre}, Corps locaux, Hermann, Paris (1962).
\end{thebibliography}
